\chapter{Introduction}
\label{chap:introduction}

\section{Problem Statement}

\subsection{Context}

Personal finance management requires users to record transactions
consistently, categorize them correctly, and make sense of their
income--expense history. Financial literacy is directly linked to the
ability to plan and to make sound long-term decisions
\cite{lusardi2014}. Traditional bookkeeping tools, however, usually
demand many manual steps: choosing the transaction type, entering the
amount, and selecting a category, a wallet, and a date. When data entry
becomes a burden, records are easily left incomplete or mislabeled, and
every downstream report becomes less trustworthy.

Natural input such as ``ph\d{o} for breakfast today, 45k paid by Momo''
reduces the number of steps, but it introduces a new engineering
problem. The system must understand Vietnamese currency units, relative
dates, several transactions in a single sentence, approximate wallet
names, and the context needed for categorization. Receipt images and
voice add further OCR/STT noise. If a speculative result is written
straight into the database, a small model error can corrupt the
balance, the budget, and every insight computed afterwards.

\subsection{The Analysis and Interpretation Problem}

Charts and filters answer well-defined questions such as ``how much did
I spend this month.'' The research value of PERFIN does not lie in
wrapping an SQL query inside a chatbot. The system must instead detect
phenomena that a single number cannot reveal: a gradually rising trend,
an abnormal spending day, the rate at which cash is being depleted, a
small recurring charge, or a goal whose progress has drifted off track.

These phenomena must be computed by algorithms that can be tested. An
LLM can understand a question and interpret a result, but it is not a
source of truth for financial computation. Delegating the entire
analysis to an LLM makes results hard to reproduce and increases the
risk of numbers that do not exist in the source data.

\subsection{Research Questions}

The project focuses on three questions:

\begin{enumerate}
  \item How can text, image, and voice input be turned into structured
  transactions while still protecting data correctness?
  \item Which deterministic algorithms are suitable for producing
  explainable insights from a personal financial history?
  \item Where should the LLM participate so that it increases
  naturalness and personalization without seizing control of computation
  or silently altering data?
\end{enumerate}

\widereportfigure{01-system-context.pdf}{System context and scope of
PERFIN. The user interacts with the mobile application; the LLM, OCR,
and STT providers are external supporting services only.}{fig:system-context}

Figure~\ref{fig:system-context} positions PERFIN as a data- and
algorithm-management system. Banking integration, shared wallets,
production authentication, and high-availability infrastructure are out
of scope for this basic-level project.

\section{Objectives}

The overall objective is to build and evaluate a personal finance
prototype in which structured data and deterministic algorithms form the
foundation, while the LLM acts as a controlled layer for language
understanding and interpretation. The specific objectives in
Table~\ref{tab:objectives} are written to be specific, measurable, and
bounded, each paired with the evidence and the acceptance threshold that
would let it be marked complete.

\begin{table}[H]
  \centering
  \caption{Objectives and evaluation criteria of the project}
  \label{tab:objectives}
  \begin{tabularx}{\textwidth}{@{}L{1.1cm}Y Y@{}}
    \toprule
    \textbf{ID} & \textbf{Specific objective} & \textbf{Evidence and acceptance criteria} \\
    \midrule
    O1 & Build a financial data model with keys, constraints, indexes, and atomic transactions. & Migrations, ERD, and rollback tests; every critical update flow must either complete fully or leave the data unchanged. \\
    O2 & Build a multimodal extraction pipeline with clarification and confirmation. & A labeled dataset and tool-call logs; no database write when a required field is missing. \\
    O3 & Implement trend, anomaly, runway, recurring, correlation, budget, and goal algorithms. & Formula tests over ordinary data, boundary data, and hand-computed expected values. \\
    O4 & Restrict the LLM to intent understanding, tool calling, and fact interpretation. & Tool schemas, numeric-faithfulness checks, and fallback tests; no number introduced beyond the supplied facts. \\
    O5 & Evaluate the operational readiness of the prototype. & Integration tests, p50/p95 latency, provider errors, and UAT, each with reproducible logs before publication. \\
    \bottomrule
  \end{tabularx}
\end{table}

The thresholds in Table~\ref{tab:objectives} are \textbf{acceptance
targets}, not results that have automatically been achieved. The report
applies the label ``measured'' only when a reproducible command,
environment, and output exist.

\section{Scope}

\subsection{Scope of Work}

\begin{table}[H]
  \centering
  \caption{Scope of work and out-of-scope items}
  \label{tab:scope}
  \begin{tabularx}{\textwidth}{@{}Y Y@{}}
    \toprule
    \textbf{In scope} & \textbf{Out of scope} \\
    \midrule
    Income, expense, wallet-transfer, and investment profit/loss transactions at the prototype level. & Open Banking integration and real bank synchronization. \\
    Text, receipt-image, and voice input; preview and confirmation. & Training new foundation, OCR, or STT models. \\
    Categories, correction feedback, suggestions, and confirmed re-tagging. & Shared wallets, group debt splitting, and real-time collaboration. \\
    Budgets, recurring bills, insights, and financial goals. & Investment advice, credit scoring, or decisions made on the user's behalf. \\
    A single \code{default_user} demo profile and a data-upgrade path. & Registration, JWT, authorization, and production multi-user isolation. \\
    Redis, worker, and fallbacks for demo and testing. & High availability, centralized monitoring, and uptime commitments. \\
    \bottomrule
  \end{tabularx}
\end{table}

The focus of the project is the data model, the correctness of data
transformations, and the analytical algorithms. The mobile interface
serves to capture input and to illustrate usability; exhaustive CRUD for
every object is not the primary measure of academic contribution.

\subsection{Assumptions and Limitations}

\begin{itemize}
  \item The default currency is VND; monetary data uses a decimal type
  in PostgreSQL. Numbers generated by the LLM must not be used as a
  source of truth.
  \item Statistical conclusions are meaningful only when the number of
  observations is large enough. The system must lower its confidence or
  decline to conclude when the data is too sparse.
  \item Personas change only the manner of expression, not the facts,
  the computation, the access rights, or the confirmation conditions.
  \item The prototype does not replace a financial professional;
  suggestions only help the user organize and understand personal data.
\end{itemize}

\section{Contributions}

The project targets five verifiable contributions:

\begin{enumerate}
  \item a constrained data model with atomic update flows;
  \item a pipeline that combines an LLM, a local parser, fuzzy matching,
  and a human in the loop;
  \item a set of deterministic analytical algorithms, separated from the
  language-generation layer;
  \item a grounded-narration mechanism in which the LLM interprets only
  facts that have already been computed;
  \item an evaluation protocol that separates algorithmic correctness,
  AI quality, performance, and the data-entry experience.
\end{enumerate}

\section{Report Structure}

Chapter~\ref{chap:introduction} presents the problem, objectives, and
scope. Chapter~\ref{chap:theory} builds the theoretical foundation on
financial data, LLMs, OCR/STT, and analytical algorithms.
Chapter~\ref{chap:results} specifies the requirements, the architectural
design, the data model, the processing flows, and the testing.
Chapter~\ref{chap:conclusion} maps objectives to evidence, states the
limitations, and proposes future work.
